\section{Research Thrusts}

% This \section header introduces all thrusts; thrusts 2 and 3 continue with \subsection.

\subsection{Thrust 1: \thrustone}
\label{sec:thrust1}

\noindent\textbf{Problem Statement.}
% TODO(glitching): 2-4 sentences. What is the gap?
%   - Fault-injection defenses are reasoned about at the source/ISA level, but the binary that
%     actually runs is produced by the compiler.
%   - Optimization, instruction selection, register allocation, and scheduling can collapse
%     redundant checks, create single points of failure, and otherwise enlarge the surface a
%     glitch can exploit -- invisibly to the developer.
%   - This compiler-introduced surface is neither characterized nor quantified today.

\smallskip
\noindent\textbf{Research Questions.}
\begin{itemize}[leftmargin=*,nosep]
  \item \textbf{RQ1.1:} Which compiler transformations (optimization passes, instruction selection, register allocation, scheduling) introduce or enlarge the glitching attack surface, and by how much?
  \item \textbf{RQ1.2:} How can the glitching attack surface of compiled code be \emph{quantified} so it can be compared across compilers, optimization levels, and target ISAs?
  \item \textbf{RQ1.3:} Do some compiler decisions \emph{reduce} the surface, revealing low-cost levers for Thrust~3?
\end{itemize}

\smallskip
\noindent\textbf{Approach Overview.}
% TODO(glitching): one paragraph tying the tasks below to the RQs above.

\subsubsection{\researchoneone}
\noindent\textbf{Goal.} % TODO(glitching)
\noindent\textbf{Approach.} % TODO(glitching): differential analysis across passes / -O levels; map source-level checks to emitted instructions.
\noindent\textbf{Expected Outcome.} % TODO(glitching)

\subsubsection{\researchonetwo}
\noindent\textbf{Goal.} % TODO(glitching)
\noindent\textbf{Approach.} % TODO(glitching): define a surface metric (e.g., exploitable single-fault sites per function) and a measurement methodology.
\noindent\textbf{Expected Outcome.} % TODO(glitching)
